%---------------
%╔═╗╔═╗╔╦╗╦ ╦╔═╗
%╚═╗║╣  ║ ║ ║╠═╝
%╚═╝╚═╝ ╩ ╚═╝╩  
%---------------
\documentclass[12pt,oneside,a4paper]{report}

% DOCUMENT SETUP
\usepackage[left=3cm, 
			right=2.5cm, 
			top=2.5cm, 
			bottom=2.5cm, 
			includehead, 
			includefoot]{geometry}

% line spacing
\usepackage{setspace}
\setstretch{1,25} % 15/12 --> 1.25

%de­fines Adobe Times Ro­man as de­fault text font
\usepackage{mathptmx}
\usepackage{times} % needed for acronym package

%PDF linking package
\usepackage[hidelinks]{hyperref}

% Language Setup
\usepackage[ngerman]{babel}
% language specific bibliography style
\usepackage[numbers]{natbib}
\usepackage[fixlanguage]{babelbib}
\selectbiblanguage{german}
% bliographystyle setup
% default style names: apalike alphadin ieeetr IEEEtranSN apalike2 alphadin 
% babel specific: babplain, babplai3, babalpha, babunsrt, bababbrv, bababbr3 unsrt 
\bibliographystyle{unsrturl}

% encoding setup
% T1 font encoding for languages that use a latin alphabet
\usepackage[T1]{fontenc} 

% enhanced input encoding handling - utf8 for äÄüÜöÖß...
\usepackage[utf8]{inputenc}
%\usepackage{ucs}%utf8x suppart

% after babel - set chapter string
\AtBeginDocument{\renewcommand{\chaptername}{}}

% enumeration
\usepackage{enumitem}
% tabular extension tabularx
\usepackage{tabularx}

% math packages
\usepackage{amsmath}
\usepackage{nicefrac}
\usepackage{amsthm}
\usepackage{amsbsy}
\usepackage{amssymb}
\usepackage{amsfonts}
\usepackage{MnSymbol}

% patches for latex
\usepackage{fixltx2e}

%special characters
\usepackage{amssymb}
\usepackage{upgreek,textgreek}

% acronym package
\usepackage[printonlyused, footnote]{acronym}

% breakable text in \seqsplit{}
\usepackage{seqsplit}

% \textmu
\usepackage{textcomp}

% package provides a way to compile sections of a document using the same preamble as the main document
\usepackage{subfiles}

% driver-independent color extension - used by listings,tabularx
\usepackage[usenames,dvipsnames,table,xcdraw]{xcolor}

% -- SYNTAX HIGHLIGHTING --
\usepackage{listings}
%% bash command line Syntax Highlighting
\lstdefinestyle{BASH_CMD}{ 
  columns=fullflexible,            % copy pasteable listings
  language=bash,
  basicstyle=\small\sffamily,
  basicstyle   = \small \ttfamily,
  keywordstyle = [1]\small \ttfamily,
  keywordstyle = [2]\small \ttfamily,
  commentstyle = \small \ttfamily,
  numbers=none,
  captionpos=b, 
  breaklines=true,
  numberstyle=\tiny,
  numbersep=3pt,
  frame=tlrb,
  columns=fullflexible,
  backgroundcolor=\color{white!20},
  linewidth=\linewidth,
  literate=                        % replace in code
     {Ö}{{\"O}}1
     {Ä}{{\"A}}1
     {Ü}{{\"U}}1
     {ß}{{\ss}}2
     {ü}{{\"u}}1
     {ä}{{\"a}}1
     {ö}{{\"o}}1
}
 % adds style BASH_CMD
%% Matlab Syntax Highlighting
\colorlet{keyword}{blue!100!black!80}
\colorlet{STD}{Lavender}
\colorlet{comment}{green!90!black!90}
\definecolor{mygreen}{rgb}{0,0.6,0}
\definecolor{mygray}{rgb}{0.5,0.5,0.5}
\definecolor{mymauve}{rgb}{0.58,0,0.82}


\lstdefinestyle{BASH_SCRIPT}{ 
  language     = bash,
  basicstyle   = \footnotesize \ttfamily,
  keywordstyle = [1]\color{keyword}\bfseries,
  keywordstyle = [2]\color{STD}\bfseries,
  commentstyle = \color{mygreen}\itshape,
  backgroundcolor=\color{white},   % choose the background color; you must add \usepackage{color} 
  columns=fullflexible,            % copy pasteable listings
                                   % or \usepackage{xcolor}
  basicstyle=\footnotesize,        % the size of the fonts that are used for the code
  breakatwhitespace=false,         % sets if automatic breaks should only happen at whitespace
  breaklines=true,                 % sets automatic line breaking
  captionpos=b,                    % sets the caption-position to bottom
  extendedchars=true,              % lets you use non-ASCII characters; for 8-bits encodings only,
                                   % does not work with UTF-8
  frame=single,                    % adds a frame around the code
  keepspaces=true,                 % keeps spaces in text, useful for keeping indentation of code
                                   % (possibly needs columns=flexible)
  numbers=left,                    % where to put the line-numbers; possible values are 
                                   % (none, left, right)
  numbersep=5pt,                   % how far the line-numbers are from the code
  numberstyle=\tiny\color{mygray}, % the style that is used for the line-numbers
  rulecolor=\color{black},         % if not set, the frame-color may be changed on line-breaks
                                   % within not-black text (e.g. comments (green here))
  showspaces=false,                % show spaces everywhere adding particular underscores; it
  	                               % overrides 'showstringspaces'
  showstringspaces=false,          % underline spaces within strings only
  showtabs=false,                  % show tabs within strings adding particular underscores
  stepnumber=1,                    % the step between two line-numbers. If it's 1, each line 
                                   % will be numbered
  stringstyle=\color{mymauve},     % string literal style
  tabsize=2,                       % sets default tabsize to 2 spaces
  title=\lstname,                  % set title name
  literate=                        % replace in code
     {Ö}{{\"O}}1
     {Ä}{{\"A}}1
     {Ü}{{\"U}}1
     {ß}{{\ss}}2
     {ü}{{\"u}}1
     {ä}{{\"a}}1
     {ö}{{\"o}}1
} % adds style BASH_SCRIPT
% Matlab Syntax Highlighting
\colorlet{keyword}{blue!100!black!80}
\colorlet{STD}{red}
\colorlet{comment}{green!90!black!90}
\definecolor{mygreen}{rgb}{0,0.6,0}
\definecolor{mygray}{rgb}{0.5,0.5,0.5}
\definecolor{mymauve}{rgb}{0.58,0,0.82}


\lstdefinestyle{LATEX}{ 
  language     = [LaTeX]{TeX},
  basicstyle   = \footnotesize \ttfamily,
  keywordstyle = [1]\color{keyword}\bfseries,
  keywordstyle = [2]\color{comment}\bfseries,
  commentstyle = \color{mygray}\itshape,
  %backgroundcolor=\color{white},   % choose the background color; you must add \usepackage{color} 
                                   % or \usepackage{xcolor}
  basicstyle=\footnotesize,        		   % the size of the fonts that are used for the code
  breakatwhitespace=false,         % sets if automatic breaks should only happen at whitespace
  columns=fullflexible,            % copy pasteable listings
  breaklines=true,                 % sets automatic line breaking
  captionpos=c,                    % sets the caption-position to bottom
  extendedchars=true,              % lets you use non-ASCII characters; for 8-bits encodings only,
                                   % does not work with UTF-8
  frame=single,                    % adds a frame around the code
  keepspaces=true,                 % keeps spaces in text, useful for keeping indentation of code
                                   % (possibly needs columns=flexible)
  numbers=left,                    % where to put the line-numbers; possible values are 
                                   % (none, left, right)
  numbersep=4pt,                   % how far the line-numbers are from the code
  numberstyle=\tiny\color{mygray}, % the style that is used for the line-numbers
  rulecolor=\color{black},         % if not set, the frame-color may be changed on line-breaks
                                   % within not-black text (e.g. comments (green here))
  showspaces=false,                % show spaces everywhere adding particular underscores; it
  	                               % overrides 'showstringspaces'
  showstringspaces=false,          % underline spaces within strings only
  showtabs=false,                  % show tabs within strings adding particular underscores
  stepnumber=1,                    % the step between two line-numbers. If it's 1, each line 
                                   % will be numbered
  stringstyle=\color{mymauve},     % string literal style
  tabsize=2,                       % sets default tabsize to 2 spaces
  title=\lstname,                  % set title name
  literate=                        % replace in code
     {Ö}{{\"O}}1
     {Ä}{{\"A}}1
     {Ü}{{\"U}}1
     {ß}{{\ss}}2
     {ü}{{\"u}}1
     {ä}{{\"a}}1
     {ö}{{\"o}}1
} % adds style LATEX
%% Matlab Syntax Highlighting
\colorlet{keyword}{blue!100!black!80}
\colorlet{STD}{Lavender}
\colorlet{comment}{green!90!black!90}
\definecolor{mygreen}{rgb}{0,0.6,0}
\definecolor{mygray}{rgb}{0.5,0.5,0.5}
\definecolor{mymauve}{rgb}{0.58,0,0.82}


\lstdefinestyle{MATLAB}{ 
  language     = Matlab,
  basicstyle   = \footnotesize \ttfamily,
  keywordstyle = [1]\color{keyword}\bfseries,
  keywordstyle = [2]\color{STD}\bfseries,
  commentstyle = \color{mygreen}\itshape,
  backgroundcolor=\color{white},   % choose the background color; you must add \usepackage{color} 
                                   % or \usepackage{xcolor}
  basicstyle=\footnotesize,        % the size of the fonts that are used for the code
  breakatwhitespace=false,         % sets if automatic breaks should only happen at whitespace
  columns=fullflexible,            % copy pasteable listings
  breaklines=false,                % sets automatic line breaking
  captionpos=c,                    % sets the caption-position to bottom
  extendedchars=true,              % lets you use non-ASCII characters; for 8-bits encodings only,
                                   % does not work with UTF-8
  frame=single,                    % adds a frame around the code
  keepspaces=true,                 % keeps spaces in text, useful for keeping indentation of code
                                   % (possibly needs columns=flexible)
  numbers=left,                    % where to put the line-numbers; possible values are 
                                   % (none, left, right)
  numbersep=5pt,                   % how far the line-numbers are from the code
  numberstyle=\tiny\color{mygray}, % the style that is used for the line-numbers
  rulecolor=\color{black},         % if not set, the frame-color may be changed on line-breaks
                                   % within not-black text (e.g. comments (green here))
  showspaces=false,                % show spaces everywhere adding particular underscores; it
  	                               % overrides 'showstringspaces'
  showstringspaces=false,          % underline spaces within strings only
  showtabs=false,                  % show tabs within strings adding particular underscores
  stepnumber=1,                    % the step between two line-numbers. If it's 1, each line 
                                   % will be numbered
  stringstyle=\color{mymauve},     % string literal style
  tabsize=2,                       % sets default tabsize to 2 spaces
  title=\lstname,                  % set title name
  literate=                        % replace in code
     {Ö}{{\"O}}1
     {Ä}{{\"A}}1
     {Ü}{{\"U}}1
     {ß}{{\ss}}2
     {ü}{{\"u}}1
     {ä}{{\"a}}1
     {ö}{{\"o}}1
} % adds style MATLAB
% Matlab Syntax Highlighting
\colorlet{keyword}{blue!100!black!80}
\colorlet{STD}{Lavender}
\colorlet{comment}{green!90!black!90}
\definecolor{mygreen}{rgb}{0,0.6,0}
\definecolor{mygray}{rgb}{0.5,0.5,0.5}
\definecolor{mymauve}{rgb}{0.58,0,0.82}


\lstdefinestyle{PYTHON}{ 
  language     = Python,
  basicstyle   = \footnotesize \ttfamily,
  keywordstyle = [1]\color{keyword}\bfseries,
  keywordstyle = [2]\color{STD}\bfseries,
  commentstyle = \color{mygreen}\itshape,
  backgroundcolor=\color{white},   % choose the background color; you must add \usepackage{color} 
                                   % or \usepackage{xcolor}
  basicstyle=\footnotesize,        % the size of the fonts that are used for the code
  columns=fullflexible,            % copy pasteable listings
  breakatwhitespace=false,         % sets if automatic breaks should only happen at whitespace
  breaklines=false,                % sets automatic line breaking
  captionpos=c,                    % sets the caption-position to bottom
  extendedchars=true,              % lets you use non-ASCII characters; for 8-bits encodings only,
                                   % does not work with UTF-8
  frame=single,                    % adds a frame around the code
  keepspaces=true,                 % keeps spaces in text, useful for keeping indentation of code
                                   % (possibly needs columns=flexible)
  numbers=left,                    % where to put the line-numbers; possible values are 
                                   % (none, left, right)
  numbersep=5pt,                   % how far the line-numbers are from the code
  numberstyle=\tiny\color{mygray}, % the style that is used for the line-numbers
  rulecolor=\color{black},         % if not set, the frame-color may be changed on line-breaks
                                   % within not-black text (e.g. comments (green here))
  showspaces=false,                % show spaces everywhere adding particular underscores; it
  	                               % overrides 'showstringspaces'
  showstringspaces=false,          % underline spaces within strings only
  showtabs=false,                  % show tabs within strings adding particular underscores
  stepnumber=1,                    % the step between two line-numbers. If it's 1, each line 
                                   % will be numbered
  stringstyle=\color{mymauve},     % string literal style
  tabsize=2,                       % sets default tabsize to 2 spaces
  title=\lstname,                  % set title name
  literate=                        % replace in code
     {Ö}{{\"O}}1
     {Ä}{{\"A}}1
     {Ü}{{\"U}}1
     {ß}{{\ss}}2
     {ü}{{\"u}}1
     {ä}{{\"a}}1
     {ö}{{\"o}}1
} % adds style PYTHON
%% Matlab Syntax Highlighting
\colorlet{keyword}{blue!100!black!80}
\colorlet{STD}{Lavender}
\colorlet{comment}{green!90!black!90}
\definecolor{mygreen}{rgb}{0,0.6,0}
\definecolor{mygray}{rgb}{0.5,0.5,0.5}
\definecolor{mymauve}{rgb}{0.58,0,0.82}


\lstdefinestyle{CPP}{ 
  language     = C++,
  basicstyle   = \footnotesize \ttfamily,
  keywordstyle = [1]\color{keyword}\bfseries,
  keywordstyle = [2]\color{STD}\bfseries,
  commentstyle = \color{mygreen}\itshape,
  backgroundcolor=\color{white},   % choose the background color; you must add \usepackage{color} 
                                   % or \usepackage{xcolor}
  columns=fullflexible,            % copy pasteable listings
  basicstyle=\footnotesize,        % the size of the fonts that are used for the code
  breakatwhitespace=false,         % sets if automatic breaks should only happen at whitespace
  breaklines=false,                % sets automatic line breaking
  captionpos=c,                    % sets the caption-position to bottom
  extendedchars=true,              % lets you use non-ASCII characters; for 8-bits encodings only,
                                   % does not work with UTF-8
  frame=single,                    % adds a frame around the code
  keepspaces=true,                 % keeps spaces in text, useful for keeping indentation of code
                                   % (possibly needs columns=flexible)
  numbers=left,                    % where to put the line-numbers; possible values are 
                                   % (none, left, right)
  numbersep=5pt,                   % how far the line-numbers are from the code
  numberstyle=\tiny\color{mygray}, % the style that is used for the line-numbers
  rulecolor=\color{black},         % if not set, the frame-color may be changed on line-breaks
                                   % within not-black text (e.g. comments (green here))
  showspaces=false,                % show spaces everywhere adding particular underscores; it
  	                               % overrides 'showstringspaces'
  showstringspaces=false,          % underline spaces within strings only
  showtabs=false,                  % show tabs within strings adding particular underscores
  stepnumber=1,                    % the step between two line-numbers. If it's 1, each line 
                                   % will be numbered
  stringstyle=\color{mymauve},     % string literal style
  tabsize=2,                       % sets default tabsize to 2 spaces
  title=\lstname,                  % set title name
  literate=                        % replace in code
     {Ö}{{\"O}}1
     {Ä}{{\"A}}1
     {Ü}{{\"U}}1
     {ß}{{\ss}}2
     {ü}{{\"u}}1
     {ä}{{\"a}}1
     {ö}{{\"o}}1
} % adds style CPP
%% Matlab Syntax Highlighting
\colorlet{keyword}{blue!100!black!80}
\colorlet{STD}{Lavender}
\colorlet{comment}{green!90!black!90}
\definecolor{mygreen}{rgb}{0,0.6,0}
\definecolor{mygray}{rgb}{0.5,0.5,0.5}
\definecolor{mymauve}{rgb}{0.58,0,0.82}


\lstdefinestyle{C}{ 
  language     = C,
  basicstyle   = \footnotesize \ttfamily,
  keywordstyle = [1]\color{keyword}\bfseries,
  keywordstyle = [2]\color{STD}\bfseries,
  commentstyle = \color{mygreen}\itshape,
  backgroundcolor=\color{white},   % choose the background color; you must add \usepackage{color} 
  columns=fullflexible,            % copy pasteable listings
                                   % or \usepackage{xcolor}
  basicstyle=\footnotesize,        % the size of the fonts that are used for the code
  breakatwhitespace=false,         % sets if automatic breaks should only happen at whitespace
  breaklines=false,                % sets automatic line breaking
  captionpos=c,                    % sets the caption-position to bottom
  extendedchars=true,              % lets you use non-ASCII characters; for 8-bits encodings only,
                                   % does not work with UTF-8
  frame=single,                    % adds a frame around the code
  keepspaces=true,                 % keeps spaces in text, useful for keeping indentation of code
                                   % (possibly needs columns=flexible)
  numbers=left,                    % where to put the line-numbers; possible values are 
                                   % (none, left, right)
  numbersep=5pt,                   % how far the line-numbers are from the code
  numberstyle=\tiny\color{mygray}, % the style that is used for the line-numbers
  rulecolor=\color{black},         % if not set, the frame-color may be changed on line-breaks
                                   % within not-black text (e.g. comments (green here))
  showspaces=false,                % show spaces everywhere adding particular underscores; it
  	                               % overrides 'showstringspaces'
  showstringspaces=false,          % underline spaces within strings only
  showtabs=false,                  % show tabs within strings adding particular underscores
  stepnumber=1,                    % the step between two line-numbers. If it's 1, each line 
                                   % will be numbered
  stringstyle=\color{mymauve},     % string literal style
  tabsize=2,                       % sets default tabsize to 2 spaces
  title=\lstname,                  % set title name
  literate=                        % replace in code
     {Ö}{{\"O}}1
     {Ä}{{\"A}}1
     {Ü}{{\"U}}1
     {ß}{{\ss}}2
     {ü}{{\"u}}1
     {ä}{{\"a}}1
     {ö}{{\"o}}1
} % adds style C
%% JSON Syntax Highlighting
\colorlet{keyword}{blue!100!black!80}
\colorlet{STD}{Lavender}
\colorlet{comment}{green!90!black!90}
\definecolor{mygreen}{rgb}{0,0.6,0}
\definecolor{mygray}{rgb}{0.5,0.5,0.5}
\definecolor{mymauve}{rgb}{0.58,0,0.82}

\newcommand\JSONnumbervaluestyle{\color{blue}}
\newcommand\JSONstringvaluestyle{\color{red}}

\newif\ifcolonfoundonthisline

\makeatletter

\lstdefinelanguage{json}
{
  showstringspaces    = false,
  keywords            = {false,true},
  alsoletter          = 0123456789.,
  morestring          = [s]{"}{"},
  morestring          = [s]{'}{'},
  stringstyle         = \ifcolonfoundonthisline\JSONstringvaluestyle\fi,
  MoreSelectCharTable =%
    \lst@DefSaveDef{`:}\colon@json{\processColon@json},
  basicstyle          = \ttfamily,
  keywordstyle        = \ttfamily\bfseries,
}

% flip the switch if a colon is found in Pmode
\newcommand\processColon@json{
  \colon@json%
  \ifnum\lst@mode=\lst@Pmode%
    \global\colonfoundonthislinetrue%
  \fi
}

\lst@AddToHook{Output}{%
  \ifcolonfoundonthisline%
    \ifnum\lst@mode=\lst@Pmode%
      \def\lst@thestyle{\JSONnumbervaluestyle}%
    \fi
  \fi
  %override by keyword style if a keyword is detected!
  \lsthk@DetectKeywords% 
}

% reset the switch at the end of line
\lst@AddToHook{EOL}%
  {\global\colonfoundonthislinefalse}

\makeatother



\lstdefinestyle{JSON}{ 
  language     = json,
  basicstyle   = \footnotesize \ttfamily,
  keywordstyle = [1]\color{keyword}\bfseries,
  keywordstyle = [2]\color{STD}\bfseries,
  commentstyle = \color{mygreen}\itshape,
  backgroundcolor=\color{white},   % choose the background color; you must add \usepackage{color} 
                                   % or \usepackage{xcolor}
  basicstyle=\footnotesize,        % the size of the fonts that are used for the code
  columns=fullflexible,            % copy pasteable listings
  breakatwhitespace=false,         % sets if automatic breaks should only happen at whitespace
  breaklines=false,                % sets automatic line breaking
  captionpos=c,                    % sets the caption-position to bottom
  extendedchars=true,              % lets you use non-ASCII characters; for 8-bits encodings only,
                                   % does not work with UTF-8
  frame=single,                    % adds a frame around the code
  keepspaces=true,                 % keeps spaces in text, useful for keeping indentation of code
                                   % (possibly needs columns=flexible)
  numbers=left,                    % where to put the line-numbers; possible values are 
                                   % (none, left, right)
  numbersep=5pt,                   % how far the line-numbers are from the code
  numberstyle=\tiny\color{mygray}, % the style that is used for the line-numbers
  rulecolor=\color{black},         % if not set, the frame-color may be changed on line-breaks
                                   % within not-black text (e.g. comments (green here))
  showspaces=false,                % show spaces everywhere adding particular underscores; it
  	                               % overrides 'showstringspaces'
  showstringspaces=false,          % underline spaces within strings only
  showtabs=false,                  % show tabs within strings adding particular underscores
  stepnumber=1,                    % the step between two line-numbers. If it's 1, each line 
                                   % will be numbered
  stringstyle=\color{mymauve},     % string literal style
  tabsize=2,                       % sets default tabsize to 2 spaces
  title=\lstname,                  % set title name
  literate=                        % replace in code
     {Ö}{{\"O}}1
     {Ä}{{\"A}}1
     {Ü}{{\"U}}1
     {ß}{{\ss}}2
     {ü}{{\"u}}1
     {ä}{{\"a}}1
     {ö}{{\"o}}1
} % adds style JSON

% HEADLINE CFG
\usepackage{fancyhdr} % Headers and footers
\usepackage{lastpage}
\usepackage{nopageno}
\setlength{\headheight}{1.5cm}
\pagestyle{fancy} % All pages have headers and footers
\fancyhead{} % Blank out the default header
\fancyfoot{} % Blank out the default footer
\fancyhead[L]{}
\fancyhead[C]{}
\fancyhead[R]{}
\fancyfoot[L]{}
\fancyfoot[C]{\thepage}
\fancyfoot[R]{}
% override plain page style for \part, \chapter or 
% \maketitle, which implicit specifies plain page style
\fancypagestyle{plain} 
{
	\fancyhead[L]{}
	\fancyhead[C]{}
	\fancyhead[R]{}
	\fancyfoot[L]{}
	\fancyfoot[C]{\thepage}
	\fancyfoot[R]{}
}
% set list pagestyle
\fancypagestyle{lists} 
{
	\fancyhead[L]{}
	\fancyhead[C]{}
	\fancyhead[R]{}
	\fancyfoot[L]{}
	\fancyfoot[C]{\thepage}
	\fancyfoot[R]{}
}

\renewcommand{\chaptermark}[1]{\markright{#1}{}}
\renewcommand{\sectionmark}[1]{\markright{#1}{}}
\renewcommand{\headrulewidth}{0pt}
\renewcommand{\footrulewidth}{0pt}

	
\usepackage{verbatim}
\usepackage{graphicx}
\usepackage{epstopdf}

% floating prevention packages
\usepackage{float}    % used with [H] positioning parameter
\usepackage{placeins} % \FloatBarrier 

% tikz packages
\usepackage{tikz}
\usepackage{caption}
\usepackage[list=true,listformat=simple]{subcaption}
\usepackage{standalone}
\usepackage{pgfplots}
\usepackage{graphicx}
\graphicspath{{./media/}}

% include only specified tex files - uncommend unneeded
\includeonly{preface/cover,
             preface/abstract,
             preface/tableofcontents,
             preface/listoffigures,
             preface/listoftables,
             preface/lstlistoflistings,
             appendix/bibliography}

%-------------------
%╔═╗╔╦╗╦═╗╦╔╗╔╔═╗╔═╗
%╚═╗ ║ ╠╦╝║║║║║ ╦╚═╗
%╚═╝ ╩ ╩╚═╩╝╚╝╚═╝╚═╝
%-------------------
\newcommand{\strLecture}{Signale, Systeme und Sensoren}
\newcommand{\strDate}{\today}
\newcommand{\strAuthorA}{Th. Gnädig}
\newcommand{\strAuthorB}{F. Gendusa}
\newcommand{\strAuthorAEmail}{thgnaedi@htwg-konstanz.de}
\newcommand{\strAuthorBEmail}{fagendus@htwg-konstanz.de}
% Versuchsbeschreibung 
\newcommand{\strTopic}{Kalibrierung von Digitalkameras}
\newcommand{\strAbstract}{In diesem Experiment befassen wir uns mit der Kalibrierung von digitalen Kameras.
Dabei finden Methoden Einsatz, wie sie der Art nach auch in der Industrie Einsatz finden.}
% hyperref customization
\hypersetup{
	pdftitle    ={\strTopic}, % title
	pdfsubject	={\strLecture}, % subject of the document
	pdfauthor	={\strAuthorA, \strAuthorB}, % author
	pdfkeywords	={}, % list of keywords
	pdfcreator	={}, % creator of the document
	pdfproducer	={}, % producer of the document
	colorlinks=false, % false: boxed links; true: colored links
	linkcolor=red, % color of internal links (change box color with linkbordercolor)
    citecolor=green, % color of links to bibliography
    filecolor=magenta, % color of file links
    urlcolor=cyan, % color of external links
	%bookmarks=true, % show bookmarks bar?
	unicode=true, % non-Latin characters in Acrobat’s bookmarks
	pdftoolbar=true, % show Acrobat’s toolbar?
	pdfmenubar=true, % show Acrobat’s menu?
    pdffitwindow=false, % window fit to page when opened
	pdfnewwindow=true % links in new PDF window
}

%-----------------------------------------
% ╔╗ ╔═╗╔═╗╦╔╗╔  ╔╦╗╔═╗╔═╗╦ ╦╔╦╗╔═╗╔╗╔╔╦╗ 
% ╠╩╗║╣ ║ ╦║║║║   ║║║ ║║  ║ ║║║║║╣ ║║║ ║  
% ╚═╝╚═╝╚═╝╩╝╚╝  ═╩╝╚═╝╚═╝╚═╝╩ ╩╚═╝╝╚╝ ╩  
%-----------------------------------------

\begin{document}
\pagenumbering{Roman} 

%\setcounter{section}{0}

\begin{titlepage}

\vspace*{-3.5cm}

\begin{flushleft}
\hspace*{-1cm} 
\includegraphics[width=15.7cm]{preface/htwg-logo}
\end{flushleft}

\vspace{1cm}

\begin{center}
	\large{
		\textbf{\strLecture} \\[2cm]
	}
	\Huge{
		\textbf{\strTopic} \\[2cm]
	}
	\Large{
		\textbf{\strAuthorA, \strAuthorB}} \\[3cm]
	\large{
		\textbf{} \\[2.3cm]
	}
	
	\large{
		\textbf{Konstanz, \strDate}
	}
\end{center}

\end{titlepage}
\thispagestyle{empty}




\begin{center}
{\Large \textbf{Zusammenfassung (Abstract)}}
\end{center}

\bigskip

\begin{center}
	\begin{tabular}{p{2.8cm}p{5cm}p{5cm}}
		Thema: & \multicolumn{2}{p{10cm}}{\raggedright\strTopic} \\
		 & & \\
		Autoren: & \strAuthorA & \href{mailto:\strAuthorAEmail}{\strAuthorAEmail} \\
		 & \strAuthorB & \href{mailto:\strAuthorBEmail}{\strAuthorBEmail} \\
		 & & \\
		Betreuer: & Prof. Dr. Matthias O. Franz & \href{mailto:mfranz@htwg-konstanz.de}{mfranz@htwg-konstanz.de} \\
		 &  Jürgen Keppler & \href{mailto:juergen.keppler@htwg-konstanz.de}{juergen.keppler@htwg-konstanz.de} \\
		 &  Martin Miller & \href{mailto:martin.miller@htwg-konstanz.de}{martin.miller@htwg-konstanz.de} \\
	\end{tabular}
\end{center}

\bigskip

\noindent
\strAbstract

\thispagestyle{lists}



\clearpage

%
% TABLE OF CONTENTS
%
\thispagestyle{lists}
%
% TABLE OF CONTENTS
%
\tableofcontents
\thispagestyle{plain}
\newpage

%
% Abbildungsverzeichnis
%
%
% Abbildungsverzeichnis
%
\phantomsection
\addcontentsline{toc}{chapter}{Abbildungsverzeichnis}
\listoffigures
\thispagestyle{lists}
\newpage

%
% Tabellenverzeichnis
%
%
% Tabellenverzeichnis
%
\phantomsection
\addcontentsline{toc}{chapter}{Tabellenverzeichnis}
\listoftables
\thispagestyle{lists}
\newpage

%
% Listingverzeichnis
%
%
% Listingverzeichnis
%
\phantomsection
\renewcommand\lstlistingname{Listing}
\renewcommand\lstlistlistingname{Listingverzeichnis}
\lstlistoflistings
\addcontentsline{toc}{chapter}{Listingverzeichnis}
\thispagestyle{lists}
\newpage


%--------------------------
% ╔═╗╦ ╦╔═╗╔═╗╔╦╗╔═╗╦═╗╔═╗ 
% ║  ╠═╣╠═╣╠═╝ ║ ║╣ ╠╦╝╚═╗ 
% ╚═╝╩ ╩╩ ╩╩   ╩ ╚═╝╩╚═╚═╝ 
%--------------------------

\pagenumbering{arabic} 
\setcounter{page}{1}
%
% CHAPTER Einleitung
%
\chapter{Einleitung}
\label{chap:EINL}
Dieser zweite Versuch der Vorlesung Signale Systeme und Sensoren beschäftigt sich mit der Kalibrierung von digitalen Kameras. Aufgeteilt wird der Versuch in folgende vier Teile:
\begin{itemize}
\setlength\itemsep{0.05em}
\item{Aufnahme und Analyse eines Grauwertkeiles}
\item{Aufnahme eines Dunkelbildes}
\item{Aufnahme eines Weißbildes}
\item{Pixelfehler}
\end{itemize}
Wie alle Sensoren sind auch digitale Bildsensoren in der Praxis nicht völlig fehlerfrei. So unterliegen sie auch systematischen Fehlern und zeigen auch Rauschverhalten auf. Im Gegensatz zu einem herkömmlichen Sensor, der genau einen Wert zu einem Zeitpunkt misst \textit{(z.B. ein einfacher Temperatursensor)} besteht ein Bildsensor aus einer Matrix voll von Helligkeitssensoren. Dies sind die sogennanten Pixel. Aufgrund von Fertigungstolleranzen weisen nicht alle Pixel exakt das selbe Verhalten auf, so dass es selbst bei einer makelos homogenen Fläche, die abgelichtet wird, bei welcher alle Pixel den selben Wert haben müssten Differenzen zwischen den Pixeln gibt. So wird also streng genommen nicht ein Bildsensor kalibriert, sondern einzelne Helligkeitssensoren(Pixel) einer Matrix(Bildsensor). In diesem Fall handelt es sich um eine Logitech Webcam mit einer Auflösung von 640 x 480 Pixeln. Bei der Kalibrierung dieser Kamera wird sich in diesem Versuch nicht um die Korrektur fehlerbehafteter Farbwiedergaben gekümmert, sondern lediglich um die Helligkeitsinterpretation der einzelnen Pixel. So werden also alle aufgenommenen Bilder in Schwarz/Weiß umgewandelt und in dem verlustfreihen png Format abgespeichert. In folgenden Abschnitten werden Fehler, die zu Bildrauschen, bzw. Bildfehlern führen aufgeführt und teilweise versucht zu korrigieren. Alle Belichtungsparameter und die Distanz von Kamera und Objekt bei allen Aufnahmen gleich.

%
% CHAPTER Versuch 1
%
\begin{lstlisting}[style=PYTHON,frame=single,
 caption=Zum einlesen der Mundharmonika Schwingung
 captionpos=b,
 label=lst:read] 
 import numpy as np
 from TekTDS2000 import *
 
 scope = TekTDS2000()
 
 input()
 scope.saveCsv(filename='mic1000Hz.csv', ch=2)
 scope.saveCsv(filename='sig1000Hz.csv', ch=1)
 scope.plot(filename='mic1000Hz.png', ch=2)
 scope.plot(filename='sig1000Hz.png', ch=1)
 scope.getSamplingInterval()
 scope.getRecordLength()
 # ausgabe bei nachmessung:
 #TekTDS2000: Sampling interval: 2e-06 s
 #TekTDS2000: Record length: 2500
 \end{lstlisting}
 
\newpage
  \begin{lstlisting}[style=PYTHON,frame=single,
   caption=Zum umwandeln der Mundharmonika Schwingung
   captionpos=b,
   label=lst:v1d]
   import numpy as np
   import matplotlib.pyplot as plt
   
   #Periodendauer
   def periodendauer():
       input_data = np.genfromtxt('tone.csv', delimiter=',')
       my_data = input_data[:,1:2];
       
       plt.figure(figsize=(10, 5), dpi=800)
       plt.xticks(range(0,271,10))
       plt.plot(my_data[0:270], label='Signal')
       plt.plot((44, 44), (-0.008, 0.008), 'k-', color='red')      # anfang Periode
       plt.plot((153, 153), (-0.008, 0.008), 'k-', color='red')    # ende Periode
       plt.ylabel("Amplitude")
       plt.xlabel("Abtastzeitpunkt")
       plt.text(57, -0.0065, '$periodendauer = 109 * 10^{-5} s$', fontsize=14)
       plt.legend()
       plt.show()
   
   
   # Fouriertransformation
   def mach_fft(data):
       fft = np.abs(np.fft.fft(data, axis=0))
       plt.figure(figsize=(10, 5), dpi=800)
       plt.xlabel("Frequenz [$Hz$]")
       plt.ylabel("Amplitude")
       plt.plot(range(0,16000,40), fft[:400], label="Fourierdomaene")
       plt.plot((910, 910), (3, 0), 'k-', color='red',label="Grundschwingung = 910 Hz") # Grundfrequenz
       #plt.plot((1820, 1820), (3, 0), 'k-', color='red') # 1 Harmonische
       #plt.plot((2730, 2730), (3, 0), 'k-', color='red') # 2 Harmonische
       plt.legend()
       plt.show()
   \end{lstlisting}

%
% CHAPTER Versuch 2
%
%
% CHAPTER Versuch 2
%
\chapter{Aufnahme eines Dunkelbildes}
In diesem Versuch wird ein Dunkelbild zur Beseitigung des Offsets jedes Pixel erstellt.
\label{chap:Aufnahme_eines_Dunkelbildes}

\section{Fragestellung, Messprinzip, Aufbau, Messmittel}
Im zweiten Teil der Kalibrierung, werden zehn Dunkelbilder aufgenommen um den Pixel Offset zu bestimmen. Es werden mehrere Bilder benötigt um den Einfluss des thermischen Rauschens der Ausleseelektronik auf das Messergebnis möglichst gering zu halten. Das Pythonskript \ref{lst:Dunkelbild} berechnet den pixelweisen Mittelwert und speichert das Durchschnittsbild ab. Das so erstellte neue Bild, entspricht dem Mittelwert der zehn aufgenommenen Bilder und sollte jetzt weitgehendst nur noch den systematischen Fehler der jeweiligen Pixel, für die gewählte Belichtungszeit, enthalten. Der Offset in der Übertragungsfunktion des jeweiligen Pixels ist auf Fertigungstoleranzen und dem so entstandenen systematischen Fehler zurückzuführen. Das Mittelwertbild \ref{img:dunkelbild_kontrastmax} enthält jetzt nur noch minimales thermisches Rauschen und kann somit als Pixelweiser Offset für spätere Korrekturen verwendet werden.
\label{chap:VERSUCH_2_FRAGESTELLUNG}

\section{Messwerte}
\begin{figure}[H]
\centering
\includegraphics[width=150mm]{dunkelbildMean_maximiert.png}
\caption{Kontrastmaximiertes Dunkelbild}
\label{img:dunkelbild_kontrastmax}
\end{figure}
\label{chap:VERSUCH_2_MESSWERTE}

\section{Auswertung und Interpretation}
\label{chap:VERSUCH_2_AUSWERTUNG}
Bei genauerem Betrachten des kontrastmaximierten Dunkelbildes sind einige graue und weiße Pixel erkennbar. Diese sind aber lediglich durch die Kontrasmaximierung zu sehen. Dies sind also Pixel, die einen Offset haben, welche dann später noch mittels Dunkelbild korrigiert werden.
Das Dunkelbild entspricht sonst den Erwartungen, da die Mehrheit der Pixel schwarz ist $farbwert \approx 0$. Ist dieser Wert exakt null findet sich hier kein Dunkelstrom \textit{(Offset)} vor. Jedoch scheinen bei genauerer Betrachtung zu viele Pixel exakt bei Wert null zu sein. Vermutlich hat die Kamera eine integrierte Funktion, die schon versucht den Dunkelstrom zu kompensieren.
Das Dunkelbild kann jetzt zum korrigieren unseres Grauwertkeilbildes verwendet werden. Der Offset der Übertragungsfunktion eines Pixels wird durch Subtraktion des korespondierenden Pixels im Dunkelbild kompensiert. Das heißt also es wird das ganze Dunkebild vom Graustufenbild subtrahiert, um alle Pixel zu korrigieren.


%
% CHAPTER Versuch 3
%
%
% CHAPTER Versuch 3
%
\chapter{Aufnahme eines Weißbildes}
\label{chap:VERSUCH_3}
Kalibrierung der Sensitivität mit Hilfe eines Weißbildes.
\section{Fragestellung, Messprinzip, Aufbau, Messmittel}
\label{chap:VERSUCH_3_FRAGESTELLUNG}
Im zweiten Teil wurde der Offset eines einzelnen Pixel im Bild mit Hilfe eines Dunkelbildes korrigiert. In diesem Teil der Kalibrierung wird die Sensitivität der Pixel korrigiert. Die Pixel weisen in ihrer Sensitivität ein hohes lineares Verhalten auf. Allerdings ist die Empfindsamkeit der einzelnen Pixel zu schulden von Fertigungstoleranzen nicht identisch. Weiter spielt die sogenannte Vignettierung mit in diesen Fehler ein. Bei dieser handelt es sich um ein ungleichmäßiges Ausleuchten der Sensorfläche durch die Optik, so dass an den Sensorrändern weniger Licht auf den Sensor fällt als in der Mitte. Deshalb gilt es auch hier eine Pixelweise Anpassung der Sensitivität vorzunehmen, damit Fertigungstolleranzen und Vignettierung kompensiert werden. Eigentlich müsste man für jede Fokuseinstellung ein eigenes Weißbild aufnehmen. Daher ist die Weiße Fläche, die abgelichtet werden soll in gleicher Entfernung, wie das Graukeilbild angebracht. 
Bei der Aufnahme des Weißbildes muss darauf geachtet werden, dass es sich bei der Aufnahmeflache um eine möglichst homogen strukturierte handelt. Diese sollte auch gleichmäßig ausgeleuchtet sein.
Dann müsste theoretisch jeder Pixel den selben Wert haben. Die Praxis zeigt jedoch, dass diese Werte nicht identisch sind, da die einen eine höhere Sensitivität haben. Diese Sensitivität entspricht gerade der Steigung einer Geraden, welche ja als Faktor auf die gemessene Größe einwirkt. Um diesen Faktor nun zu korrigieren müssen wir ihn durch den entprechenden Pixelwert des Weisbildes dividieren, um so eine Anpassung der Sensitivität zu erlangen. Da das Bild jedoch auf unbrauchbare weise verändert werden würde, wenn man einfach die korrespondierenden Pixelwerte von Weißbild mit dem Grauwertkeilbild verrechnen würde, muss das Weißbild auf einen Mittelwert von eins normiert werden. Sonst würde der Wert eines Pixels im Weißbild beispielsweise bei 178 liegen und das zu korrigierrende Bild würde drastisch verdunkelt werden. Wird das Weißbild jedoch normiert liegt der Betrag eines Pixelwertes um Eins herum. Dies kann man sich dann als prozentuale Anpassung von der Sensitivität eines Pixels vorstellen. Die Normierung des Weißbildes erreicht man, indem jeder Pixel des Bildes durch den Mittelwert des Gesamten Bildes dividiert wird.
Das Weisbild selbst ist auch wieder von thermischen Rauschen belastet, weshalb wiederum Zehn Bilder aufgenommen werden und ein Durchschnittsbild berechnet wird. Von diesem wiederum wird das Dunkelbild abgezogen, um den Offset der einzelnen Pixel zu korrigieren und anschließend das Weißbild normiert.
Dann kann das Graukeilbild durch das Weißbild dividiert werden, wobei eine Pixelweise Division darunter verstanden wird.

\section{Messwerte}
\label{chap:VERSUCH_3_MESSWERTE}
\begin{figure}[H]
\centering
\includegraphics[width=150mm]{weissbildMean_maximiert.png}
\caption{Weißbild \textit{(Mittelwert der Zehn)} Kontrastmaximiert}
\label{img:WeissBild_KontrastMax}
\end{figure}
\section{Auswertung Interpretation}
\label{chap:VERSUCH_3_AUSWERTUNG}
Das kontrastmaximierte Weißbild \ref{img:WeissBild_KontrastMax} ist in der oberen Hälfte im Vergleich zur unteren insgesamt heller. Dies hängt damit zusammen, dass das Bild an einem Fenster aufgenommen wurde und das rein scheinende Licht diese Seite wohl stärker ausgeleuchtet hat. Trotz der höheren Helligkeit im oberen Bild ist die Vignettierung zu erkennen. Diese sieht man vor allem in den Ecken des Bildes. Die schwarzen Flecken im Weißbild sind nicht als Deadpixel zu deuten. Diese sind lediglich auf eine inhomogene Oberfläche zurückzuführen(\textit{Schmutz}.

%
% CHAPTER Versuch 4
%
%
% CHAPTER Versuch 4
%
\chapter{Versuch 4}
\label{chap:VERSUCH_4}

\section{Fragestellung, Messprinzip, Aufbau, Messmittel}
\label{chap:VERSUCH_4_FRAGESTELLUNG}

\section{Messwerte}
\label{chap:VERSUCH_4_MESSWERTE}

\section{Auswertung}
\label{chap:VERSUCH_4_AUSWERTUNG}

\section{Interpretation}
\label{chap:VERSUCH_4_INTERPRETATION}

%
% CHAPTER Anhang
%
\renewcommand\thesection{A.\arabic{section}}
\renewcommand\thesubsection{\thesection.\arabic{subsection}}

\chapter*{Anhang}
\label{chap:APPENDIX}
\addcontentsline{toc}{chapter}{Anhang}
%\setcounter{chapter}{0}
\addtocounter{chapter}{1}
\setcounter{section}{0}

\section{Quellcode}
\label{chap:APPENDIX_SOURCECODE}

\subsection{Quellcode Versuch 1}
\label{chap:APPENDIX_SOURCECODE_V1}

\subsection{Quellcode Versuch 2}
\label{chap:APPENDIX_SOURCECODE_V2}

\subsection{Quellcode Versuch 3}
\label{chap:APPENDIX_SOURCECODE_V3}

\subsection{Quellcode Versuch 4}
\label{chap:APPENDIX_SOURCECODE_V4}


\section{Messergebnisse}
\label{chap:APPENDIX_MEASUREMENT_SOURCE}

%
% Literaturverzeichnis
%
%
% Literaturverzeichnis
%
\phantomsection
\addcontentsline{toc}{chapter}{Literaturverzeichnis}
\bibliography{references}
\newpage

\end{document}
%------------------------------------
% ╔═╗╔╗╔╔╦╗  ╔╦╗╔═╗╔═╗╦ ╦╔╦╗╔═╗╔╗╔╔╦╗
% ║╣ ║║║ ║║   ║║║ ║║  ║ ║║║║║╣ ║║║ ║ 
% ╚═╝╝╚╝═╩╝  ═╩╝╚═╝╚═╝╚═╝╩ ╩╚═╝╝╚╝ ╩ 
%------------------------------------